\documentclass[11pt]{article}
\usepackage[margin=1in]{geometry}
\usepackage[T1]{fontenc}
\usepackage{mathptmx}
\usepackage{xcolor}
\usepackage{enumitem}
\usepackage{parskip}
\newcommand{\resp}[1]{\par\noindent\textbf{Response.} #1}
\newcommand{\chg}[1]{\par\noindent\textbf{Changes.} \textit{#1}}
\newcommand{\pt}[1]{\par\medskip\noindent\textbf{#1}}
\begin{document}

\begin{center}
{\Large\bfseries Revision Letter}\\[4pt]
SIBGRAPI 2026 --- Paper ID 173\\
\textit{Visual Similarity Is Not Enough: Domain-Adapted Synthetic Data for
Maritime Vessel Detection}
\end{center}

\bigskip
Dear Program Chairs and Reviewers,

We thank the three reviewers for their careful and constructive reviews. We
have addressed every point raised. Beyond textual clarifications, the revision
includes \emph{new experimental evidence} produced specifically for this
rebuttal: (i)~a multi-seed replication of the pre-training epoch ablation at
the endpoints requested by Reviewer~2; (ii)~a new similarity-vs-structure
diagnostic over the full curated pool that addresses the top-$k$ curation
question without confounds; (iii)~the zero-shot evaluation of the A$'$~seq.\
arm requested by Reviewer~3; and (iv)~formal seed-paired significance tests
replacing the informal ``non-overlapping ranges'' criterion. All new scripts
and per-seed result files have been committed to the anonymised repository
cited in the paper.

A PDF with all changes highlighted (generated with \texttt{latexdiff}) is
provided alongside this letter. Section, table, and figure numbers below refer
to the revised paper.

\section*{Reviewer 1}

\pt{R1.1 --- The title claim should be softened or scoped.}
\resp{We kept the title but scoped the claim precisely: the abstract,
Section~III, and the conclusion now state that it is \emph{global, whole-image}
visual similarity that fails to predict transferability. More importantly, the
claim is now backed by a mechanism rather than a caveat: a new diagnostic
(Section~IV-A) recomputes the curation score for all 27{,}796 curated images
and shows the similarity ranking to be nearly orthogonal to the structural
axes that cause the negative transfer (see R2.W4 below).}
\chg{Abstract; Section~III (dataset paragraph); Section~IV-A; Conclusion.}

\pt{R1.2 --- The joint-rand ablation deserves open discussion.}
\resp{Agreed. The revised Section~IV-B reports the out-of-domain comparison
openly: sea-aware random placement is not merely comparable but
\emph{consistently better} out of domain ($+1.91$~pp over A$'$~joint on SMD,
all three seeds, seed-paired $p=0.005$). We now frame anchoring as an
engineering choice rather than the mechanism: we recommend in-place anchoring
when auditability and the absence of placement heuristics matter (every
synthetic position inherits a real observation; no sea-region model or
inpainting required), and random placement when cross-domain generalisation is
the priority.}
\chg{Section~IV-B (final paragraph); Section~IV-C.}

\pt{R1.3 --- Formal statistical testing.}
\resp{All headline comparisons are now assessed with two-sided seed-paired
$t$-tests (seeds matched across arms), replacing the informal criterion. The
central claims hold: A$'$~joint vs.\ B2, $+1.00$ mAP50 in every seed,
$p=0.024$; mAP50-95 $+1.50$, $p<0.001$; Recall $+2.2$, $p=0.013$; negative
transfer of A and B, both $p\le0.005$. Equivalences are also confirmed as
non-differences (A vs.\ B, $p=0.71$; B2-long vs.\ B2, $p=1.0$; joint-rand vs.\
joint in-domain, $p=0.90$). Where a comparison does \emph{not} survive testing
we now say so explicitly (see R2.W2). A footnote explains why the paired
$t$-test is the appropriate small-sample choice (with $n{=}3$ an exact
sign-permutation test cannot attain $p<0.25$).}
\chg{Section~III-C (protocol + footnote); Sections~IV-A, IV-B, IV-C
throughout.}

\pt{R1.4 --- Is crop scale conditioned on vertical position (perspective)?}
\resp{No, and the revised text says so explicitly. Section~III-B now specifies
the full A$'$~joint-rand procedure: real vessels removed by fast-marching
inpainting (Telea, radius 3~px; new reference [25]); sea region defined by an
adaptive horizon (top of the highest real bounding box), a floor at 35\% of
the image height, and an HSV refinement that excludes vertical obstacles;
rejection sampling with $\ge$75\% of the box inside the region and pairwise
IoU~$<0.2$; crop sizes drawn from the empirical width--height distribution of
all real boxes, \emph{independently of vertical position} --- perspective is
deliberately not modelled. That the unmodelled variant matches (and out of
domain exceeds) the anchored one reinforces our interpretation that the gain
stems from crop diversity and the training regime, not geometric realism.}
\chg{Section~III-B (new paragraph).}

\pt{R1.5 --- At which resolution is the size stratification computed?}
\resp{We resolved this by computational reproduction rather than recollection:
we re-implemented the size-stratified evaluation with \texttt{pycocotools},
per seed, under both candidate conventions. The convention that reproduces the
published Table~III is COCO thresholds in \emph{native image pixels}
($1920\times1080$); the alternative ($640\times640$ canvas) diverges by an
order of magnitude. Both conventions are now declared where used: Table~I's
71.6\% small-object statistic uses the $640$ canvas (matching the network
input), Table~III uses native pixels; the captions state this. Table~III
itself was regenerated from the committed script with per-seed values (minor
changes $\le0.01$, same directions), and the per-seed artifacts are in the
repository.}
\chg{Captions of Tables~I and~III; Table~III values; repository
(\texttt{ap\_by\_size.py} + per-seed JSON).}

\section*{Reviewer 2}

\pt{W1 --- Seed-level stability; sea-region procedure.}
\resp{Seed-level orderings are now reported: A$'$~joint $>$ B2 in 3/3 seeds
(likewise mAP50-95 and Recall); A and B $<$ B2 in 3/3; A$'$~seq.\ $<$ B2 in
3/3. The sea-region and sampling procedure is fully specified (see R1.4). The
anchored-vs-random comparison is discussed openly, including the out-of-domain
result favouring the random variant (see R1.2).}
\chg{Sections~III-B, IV-A, IV-B, IV-C.}

\pt{W2 --- Statistical support, in particular for the small-object claim.}
\resp{The reviewer was right on the exact metric: with the regenerated
per-seed values, $\mathrm{AR}_\mathrm{small}$ improves in all three seeds
($0.390\rightarrow0.411$, $+2.1$~points) but does not reach significance at
$n{=}3$ ($p=0.13$). The revision reports this honestly: improvements are
consistent in direction across \emph{all} six size-stratified metrics and all
seeds, significant for medium and large objects ($p<0.05$), with the largest
point estimates --- but no significance claim --- for small objects. The
abstract, results, discussion, and conclusion were all reworded accordingly.}
\chg{Abstract; Section~IV-B; Table~III; Discussion (``Operational
relevance''); Conclusion.}

\pt{W3 --- The epoch ablation uses a single seed.}
\resp{We ran the replication the reviewer requested at the endpoints. The
10-epoch point now has three seeds: mAP50 $0.8210\pm0.0036$, below the COCO
baseline in all three seeds; the 100-epoch endpoint (three seeds,
$0.7936\pm0.0049$) is significantly below the 10-epoch point (seed-paired
$p=0.028$), confirming under replication that degradation grows with exposure.
Figure~4 was regenerated: filled markers with error bars at 0, 10, and 100
epochs; open markers explicitly mark the remaining single-seed points (20,
50). Incidentally, the multi-seed endpoint strengthens the trend: the previous
single-seed 100-epoch value (0.8006) happened to be the most optimistic of the
three seeds.}
\chg{Figure~4 (figure, caption); Section~IV-A.}

\pt{W4 --- Would a stricter top-$k$ similarity curation change the outcome?}
\resp{Rather than train on one specific top-$k$ (which would confound
curation with pre-training-set size), we answer the question for the entire
family of global-similarity criteria. We recomputed the curation score for all
27{,}796 curated images under the original protocol and profiled similarity
deciles against the structural variables. The ranking is nearly orthogonal to
them: Spearman $\rho=-0.152$ against vessel area, and object density is
constant at one vessel per image, so similarity has no selective power on that
axis at all. Even the top-5\% most similar images have a median vessel area of
29\% --- versus 0.1\% operationally --- so any stricter threshold merely climbs
the same ranking towards the same structural profile. No global-similarity
threshold can close the gap; this is now stated with the numbers in
Section~IV-A, and the former ``left to future work'' sentence was replaced
accordingly. The diagnostic script and artifacts are in the repository.}
\chg{Section~III (dataset paragraph); Section~IV-A (new passage); repository
(\texttt{clip\_decile\_profile.py} + artifacts).}

\pt{W5 --- The 50/50 joint protocol is ambiguous.}
\resp{Now fully specified: 13 variations per real training image (1{,}348
real $\rightarrow$ 17{,}524 synthetic); in the balanced joint regime the real
set is repeated $13\times$, so one A$'$~joint epoch contains 35{,}048 images
(17{,}524 real + 17{,}524 synthetic, 50/50 by construction). B2-long matches
this volume with no synthetic data.}
\chg{Section~III-B; Section~III-C (B2-long definition).}

\section*{Reviewer 3}

\pt{R3.1 --- What exactly is A$'$~seq., and why 100 epochs?}
\resp{The definition is now explicit: a two-stage pipeline --- 100 epochs of
pre-training on the synthetic set, then 300 epochs of fine-tuning on real
CITRA-3D-Real only; unlike A$'$~joint, the two sources are never mixed (and
unlike B2-long, which contains no synthetic data). Extending the synthetic
stage beyond 100 epochs is counter-indicated by the monotonic degradation of
Figure~4; the results text now makes this connection. We also qualified the
$-1.31$ figure: below the baseline in all three seeds but not significant at
$n{=}3$ ($p=0.14$).}
\chg{Section~III-C; Section~IV-B.}

\pt{R3.2 --- A$'$~seq.\ is missing from the zero-shot table.}
\resp{Added: A$'$~seq.\ reaches $0.6026\pm0.0373$ mAP50 on SMD ($+1.58$ vs.\
B2, the largest seed variance in the table, hence statistically
indistinguishable from B2). This is informative: the arm that degrades
in-domain does \emph{not} degrade out of domain, indicating that the damage of
the sequential regime is specific to in-domain fine-tuning rather than a
general corruption of features. The interpretation is now in Section~IV-C.}
\chg{Table~IV; Section~IV-C.}

\pt{R3.3 --- Notation in the experimental-setup list.}
\resp{Fixed as suggested: A$'$~joint (COCO $\rightarrow$ real+synthetic
combined 50/50 in a single run).}
\chg{Section~III-C.}

\pt{R3.4 --- Table~II is too wide.}
\resp{Converted to percentages (e.g., $84.51\pm0.33$), as suggested; the
caption declares the unit and $\Delta$ in percentage points.}
\chg{Table~II.}

\section*{Corrections made on our own initiative}

While reproducing every number for this revision from committed scripts, we
found and fixed two inaccuracies; we report them here for full transparency.

\begin{enumerate}[leftmargin=1.4em]
\item \textbf{Public-set vessel area: $\sim$54\% $\rightarrow$ 39\%
(median).} The previously reported $\sim$54\% could not be reproduced. The
value measured on the actual pre-training set (27{,}796 images, committed
script) is a median of 38.9\% (mean $\sim$41\%). Figure~1, Table~I, the
abstract, and the related-work discussion were updated, and the derived scale
ratio was corrected from $\sim$500$\times$ to $\sim$400$\times$. The paper's
argument is unaffected --- if anything the corrected measurement is
conservative, and it now has full provenance.
\item \textbf{Table~IV, B2-long: $0.5755\pm0.0131 \rightarrow
0.5706\pm0.0063$.} We discovered that the previously reported value included
one seed evaluated with an intermediate checkpoint (its training run finished
later the same day and overwrote the file). Re-evaluating all three seeds with
the final checkpoints yields the corrected value; two of the three seeds
reproduce the original numbers exactly. The derived deltas in the text were
updated ($-1.1\rightarrow-1.6$~pp vs.\ B2; A$'$~joint gap
$+2.47\rightarrow+2.96$~pp). The qualitative conclusion (mild out-of-domain
harm from real-image oversampling) is unchanged, indeed slightly
strengthened.
\end{enumerate}

Table~III was likewise regenerated from the committed evaluation script with
per-seed values (changes $\le0.01$, same directions throughout); the
repository now contains every script and per-seed artifact needed to reproduce
Tables~II--IV and Figures~1 and~4.

We believe the revision addresses all concerns and thank the reviewers again
for comments that measurably improved the paper's rigour.

\bigskip
\noindent The Authors
\end{document}
